\documentclass[tikz,border=6pt]{standalone}
\usepackage{ctex}
\usepackage{amsmath}
\usepackage{amssymb}
\usepackage{pgfplots}
\pgfplotsset{compat=1.18}
\usetikzlibrary{calc,angles,quotes,arrows.meta,positioning,decorations.markings}
\begin{document}
\begin{tikzpicture}[>=Stealth]

  %==========================================================
  % 左侧：z 平面（自变量），标一个点 z = r e^{i theta}
  %==========================================================
  \begin{scope}[xshift=0cm]
    % 坐标轴
    \draw[->,gray!70] (-1.6,0) -- (2.6,0) node[right,black] {$\mathrm{Re}$};
    \draw[->,gray!70] (0,-1.8) -- (0,2.4) node[above,black] {$\mathrm{Im}$};

    % 主辐角带 (-pi, pi] 示意：负实轴为支割线
    \draw[red!70!black,line width=1.4pt,dash pattern=on 5pt off 3pt]
      (0,0) -- (-1.6,0);
    \node[red!70!black,font=\small] at (-0.9,-0.4) {支割线};

    % 单位圆参考
    \draw[gray!40,thin] (0,0) circle (1.5);

    % 点 z（取 r=1.5, theta = 110 度示意）
    \coordinate (O) at (0,0);
    \coordinate (Z) at (110:1.5);
    \draw[blue!60!black,line width=1pt] (O) -- (Z);
    \filldraw[blue!70!black] (Z) circle (1.6pt);
    \node[blue!70!black,above right=-1pt and -2pt] at (Z) {$z=re^{i\theta}$};

    % 辐角弧
    \draw[->,teal!70!black,line width=0.9pt] (0.9,0) arc (0:110:0.9);
    \node[teal!60!black,font=\small] at (45:1.18) {$\theta$};

    % 模标注
    \node[blue!60!black,font=\small] at (110:0.78) {$r$};

    \node[font=\small\bfseries] at (0.6,-1.5) {$z$ 平面};
  \end{scope}

  %==========================================================
  % 右侧：w = log z 平面（多值），竖直方向相差 2 pi i
  %==========================================================
  \begin{scope}[xshift=7.2cm]
    \def\lnr{1.6} % 代表 ln|z| 的横坐标（实部锁定）
    \def\th{1.92} % 主辐角值（约 110 度 = 1.92 弧度）

    % 坐标轴
    \draw[->,gray!70] (-1.7,0) -- (3.3,0) node[right,black] {$\mathrm{Re}\,w$};
    \draw[->,gray!70] (0,-5.0) -- (0,6.2) node[above,black] {$\mathrm{Im}\,w$};

    % 主辐角带 (-pi, pi]：用浅色横带表示主值落点区间
    \fill[orange!16] (-1.7,{-pi}) rectangle (3.3,{pi});
    \draw[orange!70!black,thin] (-1.7,{pi}) -- (3.3,{pi})
      node[right,font=\small,orange!60!black] {$\theta=\pi$};
    \draw[orange!70!black,thin,dash pattern=on 3pt off 2pt] (-1.7,{-pi}) -- (3.3,{-pi})
      node[right,font=\small,orange!60!black] {$\theta=-\pi$};
    \node[orange!55!black,font=\small,align=center] at (-1.05,0.0)
      {主辐角带\\$(-\pi,\pi]$};

    % ln|z| 的竖直参考线：只贯穿三个取值点
    \draw[blue!45,thin,dash pattern=on 2pt off 2pt]
      (\lnr,{\th-2*pi-0.5}) -- (\lnr,{\th+2*pi+0.5});
    % ln|z| 在实轴上的刻度标注
    \draw[blue!50!black,thin] (\lnr,0.12) -- (\lnr,-0.12);
    \node[blue!55!black,font=\small,below right=1pt and -2pt] at (\lnr,-0.12) {$\ln|z|$};

    % 主值点：theta 落在 (-pi, pi]
    \filldraw[blue!70!black] (\lnr,\th) circle (2.2pt);
    \node[blue!70!black,right=4pt,font=\small,align=left] at (\lnr,\th)
      {$\mathrm{Log}\,z=\ln|z|+i\theta$\\（主值）};

    % 上一支 +2 pi i
    \filldraw[purple!70!black] (\lnr,{\th+2*pi}) circle (2.2pt);
    \node[purple!70!black,right=4pt,font=\small] at (\lnr,{\th+2*pi})
      {$\mathrm{Log}\,z+2\pi i$};

    % 下一支 -2 pi i
    \filldraw[purple!70!black] (\lnr,{\th-2*pi}) circle (2.2pt);
    \node[purple!70!black,right=4pt,font=\small] at (\lnr,{\th-2*pi})
      {$\mathrm{Log}\,z-2\pi i$};

    % 间距括号标注 2 pi i（短括号，置于点左侧，远离虚轴）
    \draw[purple!60!black,line width=0.8pt]
      (\lnr-0.32,{\th+0.18}) -- (\lnr-0.32,{\th+2*pi-0.18});
    \draw[purple!60!black,line width=0.8pt] (\lnr-0.32,{\th+0.18}) -- (\lnr-0.20,{\th+0.18});
    \draw[purple!60!black,line width=0.8pt] (\lnr-0.32,{\th+2*pi-0.18}) -- (\lnr-0.20,{\th+2*pi-0.18});
    \node[purple!60!black,left=1pt,font=\small,fill=white,inner sep=1pt]
      at (\lnr-0.32,{\th+pi}) {$2\pi i$};

    \draw[purple!60!black,line width=0.8pt]
      (\lnr-0.32,{\th-0.18}) -- (\lnr-0.32,{\th-2*pi+0.18});
    \draw[purple!60!black,line width=0.8pt] (\lnr-0.32,{\th-0.18}) -- (\lnr-0.20,{\th-0.18});
    \draw[purple!60!black,line width=0.8pt] (\lnr-0.32,{\th-2*pi+0.18}) -- (\lnr-0.20,{\th-2*pi+0.18});
    \node[purple!60!black,left=1pt,font=\small,fill=white,inner sep=1pt]
      at (\lnr-0.32,{\th-pi}) {$2\pi i$};

    \node[font=\small\bfseries] at (1.0,{\th-2*pi-1.0}) {$w=\log z$ 平面};
  \end{scope}

  % 中间映射箭头
  \draw[->,line width=1.2pt,gray!80] (2.9,0.9) .. controls (4.2,1.6) and (4.4,1.0) .. (5.7,1.5);
  \node[gray!70!black,font=\small,above] at (4.3,1.6) {$\log$};

\end{tikzpicture}
\end{document}
